% Zone „Das kennst du schon“ – aus bank/gleichungen-loesen/zone.jsonl (zusammenbau v0.1)
\einheitenkopf[zone][Kennst du schon]{Das kennst du schon}
\setcounter{aufgabe}{0}

%% TODO Titel: Ich-kann-Satz fehlt in der Bank (Platzhalter: Kettenname) – Nr. 1
%% TODO Anweisung: ein Satz über der ersten Teilaufgabe fehlt in der Bank – Nr. 1
\begin{aufgabe}{Quadratische Gleichungen}
\begin{gleichungsraster}[2]
\gl{\text{Löse: $(x - 3)(x + 5) = 0$}} &
\gl{\text{Löse: $x^2 + x - 0{,}75 = 0$}} \\
\end{gleichungsraster}
\end{aufgabe}

%% TODO Titel: Ich-kann-Satz fehlt in der Bank (Platzhalter: Kettenname) – Nr. 2
%% TODO Anweisung: ein Satz über der ersten Teilaufgabe fehlt in der Bank – Nr. 2
\begin{aufgabe}{Äquivalenzumformungen und Probe}
\begin{gleichungsraster}[2]
\gl{\text{Löse: $3x + 4 = 19$}} &
\gl{\text{Löse: $0{,}5x - 1{,}5 = 2$}} \\
\end{gleichungsraster}
\end{aufgabe}

%% TODO Titel: Ich-kann-Satz fehlt in der Bank (Platzhalter: Kettenname) – Nr. 3
%% TODO Anweisung: ein Satz über der ersten Teilaufgabe fehlt in der Bank – Nr. 3
\begin{aufgabe}{Terme ausmultiplizieren, zusammenfassen, ausklammern, binomische Formeln vorwärts und rückwärts, Polynomdivision}
\begin{teile}
\teil Klammere $x$ aus: $x^3 - 7x$
\teil Klammere $x^2$ aus: $0{,}5x^4 - 3x^2$
\end{teile}
\end{aufgabe}

%% TODO Titel: Ich-kann-Satz fehlt in der Bank (Platzhalter: Kettenname) – Nr. 4
%% TODO Anweisung: ein Satz über der ersten Teilaufgabe fehlt in der Bank – Nr. 4
\begin{aufgabe}{Potenzen, n-te Wurzeln und Logarithmus}
\begin{gleichungsraster}[2]
\gl{\text{Löse: $x^3 = 8$}} &
\gl{\text{Löse: $x^5 = -32$}} \\
\end{gleichungsraster}
\end{aufgabe}

%% TODO Titel: Ich-kann-Satz fehlt in der Bank (Platzhalter: Kettenname) – Nr. 5
%% TODO Anweisung: ein Satz über der ersten Teilaufgabe fehlt in der Bank – Nr. 5
\begin{aufgabe}{Die e-Funktion ist immer positiv und nie null, e hoch null ist eins}
\begin{teile}
\teil $f(x) = 3e^x$. Wie groß ist $f(0)$? f(0) = \leerfeld
\teil $g(x) = 2e^x - 5$. Wie groß ist $g(0)$? g(0) = \leerfeld
\end{teile}
\end{aufgabe}

%% TODO Titel: Ich-kann-Satz fehlt in der Bank (Platzhalter: Kettenname) – Nr. 6
%% TODO Anweisung: ein Satz über der ersten Teilaufgabe fehlt in der Bank – Nr. 6
\begin{aufgabe}{Ableitung einer ganzrationalen Funktion und eines Produkts mit e-Funktion bilden}
\begin{teile}
\teil $f(x) = x^3 - 6x^2$. Bilde $f'(x)$ und berechne $f'(1)$. f'(x) = \leerfeld , f'(1) = \leerfeld
\teil $f(x) = -0{,}25x^4 + 2{,}5x^2$. Bilde $f'(x)$ und berechne $f'(2)$. f'(x) = \leerfeld , f'(2) = \leerfeld
\end{teile}
\end{aufgabe}

%% TODO Titel: Ich-kann-Satz fehlt in der Bank (Platzhalter: Kettenname) – Nr. 7
%% TODO Anweisung: ein Satz über der ersten Teilaufgabe fehlt in der Bank – Nr. 7
\begin{aufgabe}{Sinus am Einheitskreis}
\begin{teile}
\teil Löse $\mathrm{sin}(x) = 1$ für $0 \le x < 2\pi$. x = \leerfeld
\teil Löse $\mathrm{sin}(x) = -1$ für $0 \le x < 2\pi$. x = \leerfeld
\end{teile}
\end{aufgabe}

%% TODO Titel: Ich-kann-Satz fehlt in der Bank (Platzhalter: Kettenname) – Nr. 8
%% TODO Anweisung: ein Satz über der ersten Teilaufgabe fehlt in der Bank – Nr. 8
\begin{aufgabe}{Punkte am Graphen ablesen und zwei Graphenpunkte mit vorgegebenem Abstand in x und y finden, Maßstab (Längeneinheit in Meter)}
\begin{teile}
\teil Lies am Graphen ab: $f(4)$.

\begin{ksys}[xmin=-1,xmax=5,ymin=-1,ymax=6,ablesen] \parabel{1}{2}{1}{f} \end{ksys}
f(4) = \leerfeld
\teil Lies am Graphen ab: zwei Punkte des Graphen auf gleicher Höhe, die 4 Einheiten waagerecht auseinander liegen.

\begin{ksys}[xmin=-1,xmax=5,ymin=-1,ymax=6,ablesen] \parabel{1}{2}{1}{f} \end{ksys}
\end{teile}
\end{aufgabe}

%% TODO Titel: Ich-kann-Satz fehlt in der Bank (Platzhalter: Kettenname) – Nr. 9
%% TODO Anweisung: ein Satz über der ersten Teilaufgabe fehlt in der Bank – Nr. 9
\begin{aufgabe}{Äquivalenzumformungen und Probe – Fehler finden}
\begin{teile}
\teil Finde den Fehler und rechne richtig. \rechnung{x^2 &= 8x &&\mid :x \\ x &= 8}
\end{teile}
\end{aufgabe}

%% TODO Titel: Ich-kann-Satz fehlt in der Bank (Platzhalter: Kettenname) – Nr. 10
%% TODO Anweisung: ein Satz über der ersten Teilaufgabe fehlt in der Bank – Nr. 10
\begin{aufgabe}{Äquivalenzumformungen und Probe – selbst rechnen}
\begin{gleichungsraster}[2]
\gl{\text{Löse: $x^3 = 5x^2$}} & \\
\end{gleichungsraster}
\end{aufgabe}
